% 12-decoder.tex — 解码器逆向路径与重建
% Source: chapters/12-decoder-inverse-path-and-reconstruction.md (expanded)

\chapter{解码器逆向路径与重建}\label{ch:decoder}

\section{本章目标}

\begin{itemize}
  \item 理解解码器的完整处理管线
  \item 掌握从码流到图像的每一步逆操作
  \item 理解解码器如何验证码流完整性
  \item 了解解码器的错误处理策略
\end{itemize}

\section{背景与标准语义}

解码器是编码器的严格逆过程。每一步编码操作都有对应的逆操作，且执行顺序完全相反：

\begin{lstlisting}[numbers=none]
编码顺序: NLT→ → MCT→ → DWT→ → precinct→ → quantize→ → pack
解码顺序: unpack → dequantize → precinct→ → DWT← → MCT← ← NLT←
\end{lstlisting}

\section{完整解码流程}

\subsection{码流解析阶段}

\begin{lstlisting}[numbers=none]
xs_dec_bitstream()
  │
  ├── xs_parse_head()          // 解析所有 header markers
  │     ├── SOC (0xFF10)
  │     ├── [CAP] (0xFF50)
  │     ├── PIH (0xFF12)       // 图像参数
  │     ├── CDT (0xFF13)       // 分量表
  │     ├── WGT (0xFF14)       // 权重表
  │     ├── [NLT] (0xFF16)
  │     ├── [CWD] (0xFF17)
  │     ├── [CTS] (0xFF18)
  │     ├── [CRG] (0xFF19)
  │     └── [COM] (0xFF15)
  │
  ├── ids_construct()          // 构建内部数据结构
  │
  └── for each slice:
        SLH (0xFF20)           // 解析 slice header
        for each precinct:
          └── unpack_precinct() // 解包 precinct
\end{lstlisting}

\subsection{Precinct 解包}

\texttt{unpack\_precinct()}（\texttt{packing.c:780}）：

\begin{lstlisting}[numbers=none]
unpack_precinct()
  │
  ├── 读取 precinct header:
  │     ├── Lprc (precinct 长度)
  │     ├── Q (量化参数)
  │     ├── R (细化参数)
  │     └── per-band method signaling
  │
  ├── compute_gtli_tables(Q, R, ...)  // 恢复 GTLI 表
  │
  └── for each subpacket:
        ├── 读取 packet header (Dr, Ldata, Lgcli, Lsign)
        ├── unpack_gclis_significance()  // 解码显著性标志
        ├── unpack_gclis()               // 解码 GCLI 值
        ├── unpack_data()                // 解码系数 bitplane
        └── [unpack_sign()]              // 解码符号（Fs=1 时）
\end{lstlisting}

\subsection{逆量化}

\begin{lstlisting}[numbers=none]
dequantize_precinct()
  for each band, each row:
    dequant(data, width, gclis, N_g, gtli[band], Qpih)
\end{lstlisting}

\subsection{逆 Precinct 到图像}

\texttt{precinct\_to\_image()}（\texttt{precinct.c:172}）将 sign-magnitude 系数转换回有符号整数，再左移 $F_q$ 位恢复精度。

\subsection{逆 DWT}

\begin{lstlisting}[numbers=none]
dwt_inverse_transform()
  for each component k:
    // Phase 2 逆: 纯水平（从最深层开始）
    for d = NLx - 1 downto NLy:
      inverse_horizontal_1d_dwt()
    // Phase 1 逆: 垂直+水平
    for d = NLy - 1 downto 0:
      inverse_horizontal_1d_dwt()
      inverse_vertical_1d_dwt()
\end{lstlisting}

\subsection{逆 MCT}

\begin{lstlisting}[numbers=none]
mct_inverse_transform()
  switch (color_transform):
    case NONE:  // 无操作
    case RCT:   mct_inverse_rct()
    case TETRIX: mct_inverse_tetrix()
\end{lstlisting}

如果是 Bayer：\texttt{mct\_4\_to\_bayer()} 将 4 分量合并回单分量。

RCT 逆变换：$G = Y - \lfloor(C_b + C_r)/4\rfloor$, $R = G + C_r$, $B = G + C_b$。

\subsection{逆 NLT}

\begin{lstlisting}[numbers=none]
nlt_inverse_transform()
  switch (Tnlt):
    case NONE:      nlt_inverse_linear()
    case QUADRATIC: nlt_inverse_quadratic()
    case EXTENDED:  nlt_inverse_extended()
\end{lstlisting}

\subsection{后处理}

\begin{lstlisting}[numbers=none]
xs_dec_postprocess_image()
  // 如果是 Tetrix（Bayer），恢复原始 Bayer 尺寸
\end{lstlisting}

\section{解码器验证}

解码器在多个层次验证码流完整性：

\subsection{Marker 验证}

\begin{itemize}
  \item 每个 marker 的 ID 必须匹配预期值
  \item 长度字段必须与内容一致
  \item PIH 中的参数必须通过 \texttt{xs\_config\_validate()} 的校验
\end{itemize}

\subsection{Precinct 验证}

\begin{lstlisting}[language=C,numbers=none]
// GCLI 长度验证
if (subpkt_len != (info->gcli_len[subpkt] << 3))
    return -1;

// DATA 长度验证
if (subpkt_len != (info->data_len[subpkt] << 3))
    return -1;

// SIGN 长度验证（如果 Fs=1）
if (subpkt_len != (info->sign_len[subpkt] << 3))
    return -1;
\end{lstlisting}

每个 subpacket 的实际解码长度必须与 header 中声明的长度一致。不一致表示码流损坏。

\subsection{EOC 验证}

码流末尾必须有 EOC marker（0xFF11）。

\section{流程图}

\begin{figure}[H]
\centering
\begin{tikzpicture}[
    node distance=0.5cm,
    block/.style={rectangle, draw, fill=blue!10, text width=3.5cm, text centered, minimum height=0.7cm, font=\small},
    arrow/.style={-{Stealth[length=2.5mm]}, thick}
]
    \node[block] (input) {输入码流};
    \node[block, below=of input] (parse) {\texttt{xs\_parse\_head}\\解析 marker};
    \node[block, below=of parse] (ids) {\texttt{ids\_construct}\\重建空间索引};
    \node[block, below=of ids] (unpack) {逐 slice 逐 precinct\\解包};
    \node[block, below=of unpack] (deq) {\texttt{dequant}\\反量化};
    \node[block, below=of deq] (prec) {\texttt{precinct\_to\_image}\\$F_q$ 逆转换};
    \node[block, below=of prec] (idwt) {\texttt{dwt\_inverse\_transform}\\逆 DWT};
    \node[block, below=of idwt] (imct) {\texttt{mct\_inverse\_transform}\\逆 MCT};
    \node[block, below=of imct] (inlt) {\texttt{nlt\_inverse\_transform}\\逆 NLT};
    \node[block, below=of inlt] (output) {输出图像};

    \draw[arrow] (input) -- (parse);
    \draw[arrow] (parse) -- (ids);
    \draw[arrow] (ids) -- (unpack);
    \draw[arrow] (unpack) -- (deq);
    \draw[arrow] (deq) -- (prec);
    \draw[arrow] (prec) -- (idwt);
    \draw[arrow] (idwt) -- (imct);
    \draw[arrow] (imct) -- (inlt);
    \draw[arrow] (inlt) -- (output);
\end{tikzpicture}
\caption{解码器处理流程}
\label{fig:decoder-flowchart}
\end{figure}

\section{实现映射}

\begin{implmap}
\begin{tabular}{@{}L{2.8cm}L{3.5cm}L{4.5cm}L{4.5cm}@{}}
\toprule
标准概念 & 入口文件 & 关键函数/结构 & 说明 \\
\midrule
解码器入口 & \btt{xs\_dec.c:227} & \btt{xs\_dec\_bitstream()} & 解码一帧 \\
Probe & \btt{xs\_dec.c:105} & \btt{xs\_dec\_probe()} & 解析 header 不解码数据 \\
Marker 解析 & \btt{xs\_markers.c:673} & \btt{xs\_parse\_head()} & 解析所有 header \\
Precinct 解包 & \btt{packing.c:780} & \btt{unpack\_precinct()} & 解包一个 precinct \\
GCLI 解包 & \btt{packing.c:537} & \btt{unpack\_gclis()} & 解码 GCLI \\
数据解包 & \btt{packing.c:201} & \btt{unpack\_data()} & 解码系数 bitplane \\
反量化 & \btt{quant.c:157} & \btt{dequant()} & bitplane $\to$ 系数值 \\
Precinct $\to$ 图像 & \btt{precinct.c:172} & \btt{precinct\_to\_image()} & 系数 $\to$ 图像（含 $F_q$ 逆转换） \\
逆 DWT & \btt{dwt.c:143} & \btt{dwt\_inverse\_transform()} & 小波逆变换 \\
逆 MCT & \btt{mct.c:358} & \btt{mct\_inverse\_transform()} & 色彩逆变换 \\
逆 NLT & \btt{nlt.c:266} & \btt{nlt\_inverse\_transform()} & NLT 逆变换 \\
\bottomrule
\end{tabular}
\end{implmap}

\begin{codefact}
本章关键结论依据以下函数：
\begin{itemize}
  \item \texttt{xs\_dec\_bitstream()} --- \texttt{xs\_dec.c} --- 解码器总入口
  \item \texttt{xs\_parse\_head()} --- \texttt{xs\_markers.c:673} --- Marker 解析
  \item \texttt{unpack\_precinct()} --- \texttt{packing.c:780} --- Precinct 解包
  \item \texttt{dequant()} --- \texttt{quant.c:157} --- 反量化
  \item \texttt{precinct\_to\_image()} --- \texttt{precinct.c:172} --- 系数到图像（含 $F_q$ 逆转换）
  \item \texttt{dwt\_inverse\_transform()} --- \texttt{dwt.c:143} --- 逆 DWT
  \item \texttt{mct\_inverse\_transform()} --- \texttt{mct.c:358} --- 逆 MCT
\end{itemize}
\end{codefact}

\section{常见误解}

\begin{misunderstanding}
\begin{enumerate}
  \item \textbf{``解码器可以跳过验证''} --- 不建议。虽然可以跳过长度检查来处理轻微损坏的码流，但这可能导致未定义行为。生产代码应始终验证。
  \item \textbf{``解码器需要和编码器完全相同的配置''} --- 不需要。解码器从码流的 marker 段读取所有配置参数。\texttt{xs\_dec\_probe()} 可以只解析 header 来确定图像尺寸和格式。
  \item \textbf{``解码器的内存需求与编码器相同''} --- 不一定。解码器不需要 rate control 的预算表，内存需求可能略低。
  \item \textbf{``逆 DWT 和正 DWT 一定完全匹配''} --- 在整数域内使用 5/3 lifting 时，是的。整数 lifting 的正逆变换在没有量化的情况下是完全可逆的。
\end{enumerate}
\end{misunderstanding}

\section{与前后章衔接}

\begin{itemize}
  \item \textbf{输入}：本章的输入来自第 \ref{ch:packing} 章（打包）输出的 JPEG XS 码流
  \item \textbf{输出}：解码后的图像数据恢复为 \texttt{xs\_image\_t} 格式，与第 \ref{ch:signal-model} 章定义的信号模型一致
  \item \textbf{逆序参照}：解码步骤严格逆向第 \ref{ch:overview} 章编码器管线中的各模块（NLT $\to$ MCT $\to$ DWT $\to$ 打包）
\end{itemize}

\section{本章小结}

解码器是编码器的严格逆过程，通过逐层逆操作（unpack $\to$ dequant $\to$ DWT$\leftarrow$ $\to$ MCT$\leftarrow$ $\to$ NLT$\leftarrow$）从码流恢复图像。码流完整性通过多层验证（marker、subpacket 长度、EOC）保证。

\begin{engineeringnote}
以下``解码器可以比编码器更简单''是基于代码结构的推断。标准本身不要求解码器更简单，但 libjxs 的实现中解码器确实不包含 \texttt{rate\_control}、\texttt{fill\_gcli\_budget\_table} 等编码端模块。具体省略哪些模块取决于实现选择。
\end{engineeringnote}

JPEG XS 的设计使得解码器可以比编码器更简单——不需要 rate control、GCLI 方法选择等复杂逻辑。

接下来的章节将讲解特殊路径（Bayer/Tetrix）和端到端编码追踪。
